\documentclass[12pt,letterpaper]{report}
\usepackage{graphicx}
\usepackage{scrextend}
\usepackage{vmargin}
\usepackage{graphicx}
\usepackage{multirow}
\usepackage[utf8]{inputenc}
\usepackage[spanish]{babel}
\usepackage{multicol}
\usepackage{enumerate}
\usepackage{float}
\usepackage{amsmath, amsthm, amssymb, amsfonts}
\usepackage[usenames]{color}
\parindent=0mm
\pagestyle{empty}
\definecolor{miorange}{rgb}{0.91, 0.43, 0.0}
\begin{document}
\setmargins{2.5cm}      
{1.5cm}                     
{2cm}  
{24cm}                    
{10pt}                          
{1cm}                          
{0pt}                             
{2cm}
\begin{flushright}
\today
\end{flushright}
\vspace{-1.75cm}
\section*{Tarea Clase 3}
\subsection*{Efectúa las siguientes operaciones:}
\begin{enumerate}
\item \begin{equation*}
5x^3(5x^2+2x+3)    
\end{equation*}
\item \begin{equation*}
    (3x^2-6x+7)(4ax^2)
\end{equation*}
\item \begin{equation*}
    (a+b)(a^3+3a^2b+3ab^2+b^3)
\end{equation*}
\item \begin{equation*}
    4a^3(10a^2+2a-5)
\end{equation*}
\item \begin{equation*}
    (5x+2)^2
\end{equation*}
\item \begin{equation*}
    (6x+3)(6x-3)
\end{equation*}
\item \begin{equation*}
    (x+5)(x^5+2x^3+6x^2+4x+9)
\end{equation*}
\end{enumerate}
Las respuestas de las tareas pueden ser enviadas al siguiente correo giovannilopez9808@gmail.com con el asunto entre corchetes de la siguiente manera [Tarea #] y el nombre del archivo con sus apellidos y nombre, ya sea el nombre completo o solo una parte, por ejemplo \textit{LopezPadilla.pdf}, \textit{LopezGiovanni.pdf}, \textit{LopezPadillaGiovanniGamaliel.pdf}. Las respuestas envienlas en formato pdf, estas pueden ser totalmente digitales (que escriban todo en la computadora) o que sean fotografías, pero por favor, envielo en solo un documento.
\end{document}